%% ===================== MERGED DISCUSSION + DECLARATIONS =====================
\section{Discussion}
\label{sec:discussion}

\subsection*{What is new here}
The individual ingredients are classical: quadratic reciprocity, the quadratic-residue
obstruction to primitive roots, the consecutive-prime bias of~\cite{LOS2016}, and the
Kummer-degree density computations of~\cite{Matthews1976,MoreeStevenhagen2014}. The
contributions of this paper are: (i)~the observation that these combine into deterministic
exclusion laws on both axes --- a gap law for consecutive primes at base $10$
(Theorem~\ref{thm:exclusion}) and a same-prime pair law with a triple corollary
(Theorem~\ref{thm:pair}, Theorem~\ref{thm:triple}). The gap law is the base-$10$ instance of
the phenomenon identified by Tinkov\'a--Waxman--Zindulka~\cite{TinkovaWaxmanZindulka2023},
and we claim no priority for it; (ii)~a measurement, at the scale of $10^9$ primes, of the
consecutive-prime Artin correlation, its gap profile, and the $\omega$ repulsion, together
with the cross-base correlations $\phi(a,b)$ and their joint-density comparison; (iii)~a
descriptive reduction showing that, depending on the weighting convention, $96.1$--$99.2\%$
(modulus $120$) and $97.8$--$99.7\%$ (modulus $840$) of the consecutive-prime anticorrelation
lies between joint-residue cells, and that $95\%$ of the mean
cross-base correlation is removed by conditioning on the fine signature of $p-1$ --- that is,
both correlations are largely accounted for by the same elementary structure
--- the residues of $p$ and the factorisation of $p-1$ --- rather than by an interaction
(reductions of different estimands on different selected populations, not a complete
decomposition); and (iv)~a measured audit of the computational content
of those correlations, which finds no computational use for them within the estimators and the
workload tested (\S\ref{sec:content}). The exclusion laws are immediate
applications of standard quadratic reciprocity; their simplicity is a feature, not a
limitation, and the same is true of the audit.

\subsection*{Relation to runs of Artin primes}
Pollack~\cite{Pollack2014} proved (under GRH) that for every nonsquare $g\neq-1$ and every
$m$ there are infinitely many runs of $m$ consecutive primes each having $g$ as a primitive
root. Baker and Pollack~\cite{BakerPollack2016} proved unconditionally that if $Q$ contains
at least $\exp(Cm)$ multiplicatively independent primes, then infinitely many runs of $m$
consecutive primes each have some element of $Q$ as a primitive root, with the primes lying
in a bounded-length interval. Our results add an arithmetic constraint to the base-specific
version: for base $10$, no run of all-Artin consecutive primes can cross a gap
$\equiv 20 \pmod{40}$. We have not inferred longer-run frequencies from the pair statistic,
and we emphasise that the existence results above concern rare favourable configurations
rather than the typical behaviour measured here.

\subsection*{Relation to the Kummer-degree theory}
The pointwise obstruction of Theorem~\ref{thm:pair} and the Kummer-degree heuristic of
\S\ref{sec:model} are both governed by the same quadratic class $\sqf(ab)$: in the first it
appears as a congruence barring half the primes, in the second as a member of the group of
quadratic classes that can reduce a degree. We do \emph{not} claim these are the same
phenomenon. The product class belongs to that group for every pair, whereas it lowers a
degree only when its conductor divides the relevant cyclotomic level, and
\S\ref{sec:model} exhibits pairs separating the two. The agreement of the
Matthews--Moree--Stevenhagen densities with measurement to $0.012\%$ mean relative error is
a check of the density theory at this scale, not a new derivation of it.

\subsection*{Future work}
Four extensions suggest themselves. (1)~\emph{Other bases}: the same pipeline with
$a=2,3,5,6,7$ would test the conductor prediction after Corollary~\ref{cor:coupling}
(conductor $8$ for $a=2$, so exclusion classes mod $8$; conductor $24$ for $a=6$).
(2)~\emph{Scale}: extending the consecutive census to $10^{12}$ would help distinguish decay
from persistence in $\delta(x)$. (3)~\emph{Theory}: a derivation of $\delta(g)$ from the
singular series of~\cite{HardyLittlewood1923} combined with Hooley's method, at least for
the dominant channels, would replace empirical conditioning with a formula. (4)~\emph{Audit
protocol}: the measurement in \S\ref{sec:content} is a general recipe --- compare a claimed
predictive advantage against the elementary necessary-condition baseline, using a proper
scoring rule, and report end-to-end cost including feature acquisition --- and it would be
worth applying to other claims of algorithmic value from arithmetic correlations.

\section*{Acknowledgements}
The author thanks the developers of the open-source tools used in this work
(GMP, Python, NumPy, SymPy, Matplotlib).

\subsection*{Use of AI assistance}
Large-language-model AI assistants were used in code development, literature-search
assistance, drafting and editing, and initial formulation of the elementary proofs of
Theorems~\ref{thm:exclusion} and~\ref{thm:pair}. The computational audit of
\S\ref{sec:content} was designed with AI assistance and independently critiqued by a
separate AI system before the measurements were made. All numerical outputs were produced by
executing deterministic programs; no number in this paper was synthesised by a language
model. The author takes responsibility for the manuscript.

\section*{Declaration of interest}
No potential conflict of interest was reported by the author.

\section*{Funding}
This research received no specific grant from any funding agency in the public, commercial,
or not-for-profit sectors.

\section*{Data availability}
All code required to reproduce every number in this paper is supplied in the accompanying
reproduction package: the C segmented sieve and census programs, the independent
recomputation, the analysis scripts, the verification driver for the tables of
\S\ref{sec:content}, and the Lean development for the exclusion laws. The code, manuscript,
results and Lean development are at
\url{https://github.com/jpbald93/artin-correlations}. The per-prime dataset itself
($1.8$~GB CSV covering all $50{,}847{,}531$ primes from $7$ to $999{,}999{,}937$: gap,
$\omega(p-1)$, Artin indicator, and small residues) is not deposited because of its size,
but it is exactly and deterministically reproducible from the included C source on a single
machine, and is available from the author on request.
